\chapter{Classical Cauchy ($\alpha = 1$) and Gaussian ($\alpha = 2$) Heat Kernels}

This chapter mirrors \texttt{FracLap/WholeLine/ClassicalKernels.lean}, constructing concrete fractional heat kernel instances for the classical cases $\alpha = 1$ (Cauchy/Poisson kernel) and $\alpha = 2$ (Gaussian heat kernel) and proving both via direct Fourier inversion of the semigroup symbol $e^{-t\lvert \xi \rvert^\alpha}$.

\begin{definition}\label{def:cauchy_heat_kernel}
\lean{FracLap.cauchyHeatKernel}
\uses{def:fractional_heat_kernel}
\leanok
At fractional order $\alpha = 1$, the \emph{Cauchy (Poisson) heat kernel} is
\[
  P_1(x, t) = \frac{t}{\pi (t^2 + x^2)} \qquad (x \in \mathbb{R},\ t > 0),
\]
equipped with algebraic decay constant $C = \frac{2}{\pi} > 0$, satisfying evenness, self-similar scaling $P_1(x, t) = t^{-1} P_1(x / t, 1)$, nonnegativity, radial monotonicity, the decay bound $P_1(x, 1) \le \frac{2}{\pi (1 + \lvert x \rvert)^2}$, spatial integrability, and unit mass $\int_{\mathbb{R}} P_1(x, t)\, \mathrm{d}x = 1$.
\end{definition}

\begin{definition}\label{def:gaussian_heat_kernel}
\lean{FracLap.gaussianHeatKernel}
\uses{def:fractional_heat_kernel}
\leanok
At classical order $\alpha = 2$, the \emph{Gaussian heat kernel} is
\[
  P_2(x, t) = \frac{1}{\sqrt{4 \pi t}} \exp\Bigl(-\frac{x^2}{4t}\Bigr) \qquad (x \in \mathbb{R},\ t > 0),
\]
equipped with algebraic decay constant $C = \frac{32}{\sqrt{4\pi}} > 0$, satisfying evenness, self-similar scaling $P_2(x, t) = t^{-1/2} P_2(x t^{-1/2}, 1)$, nonnegativity, radial monotonicity, the Taylor cubic decay bound $P_2(x, 1) \le \frac{32}{\sqrt{4\pi}(1 + \lvert x \rvert)^3}$, Gaussian integrability, and unit mass $\int_{\mathbb{R}} P_2(x, t)\, \mathrm{d}x = 1$.
\end{definition}

\begin{definition}\label{def:fourier_inv_symbol}
\lean{FracLap.fourierInvFracSemigroupSymbol}
\leanok
For $\alpha, t, x \in \mathbb{R}$, the inverse Fourier transform of the fractional heat semigroup symbol $\xi \mapsto e^{-t \lvert \xi \rvert^\alpha}$ is defined by
\[
  \mathcal{F}^{-1}\bigl(e^{-t \lvert \cdot \rvert^\alpha}\bigr)(x) = \frac{1}{2\pi} \int_{\mathbb{R}} e^{\mathrm{i} x \xi}\, e^{-t \lvert \xi \rvert^\alpha}\, \mathrm{d}\xi.
\]
\end{definition}

\begin{lemma}\label{lem:integral_ioi_cexp_neg_mul}
\lean{FracLap.integral_Ioi_cexp_neg_mul}
\leanok
For every complex parameter $z \in \mathbb{C}$ with $\operatorname{Re}(z) > 0$,
\[
  \int_0^\infty e^{-z \xi}\, \mathrm{d}\xi = \frac{1}{z}.
\]
\end{lemma}

\begin{proof}
\leanok
The integrand is dominated on $(0, \infty)$ by $e^{-\operatorname{Re}(z)\xi} \in L^1(0, \infty)$. An antiderivative is $F(\xi) = (-z)^{-1} e^{-z\xi}$, which satisfies $\lim_{\xi \to +\infty} F(\xi) = 0$ since $\lvert e^{-z\xi} \rvert = e^{-\operatorname{Re}(z)\xi} \to 0$, so the improper fundamental theorem of calculus yields $\int_0^\infty e^{-z\xi}\, \mathrm{d}\xi = 0 - F(0) = z^{-1}$.
\end{proof}

\begin{theorem}\label{thm:fourier_inv_cauchy}
\lean{FracLap.fourierInvFracSemigroupSymbol_one}
\uses{def:cauchy_heat_kernel, def:fourier_inv_symbol, lem:integral_ioi_cexp_neg_mul}
\leanok
For every $t > 0$ and $x \in \mathbb{R}$, the inverse Fourier transform of $\xi \mapsto e^{-t \lvert \xi \rvert}$ equals the Cauchy kernel:
\[
  \frac{1}{2\pi} \int_{\mathbb{R}} e^{\mathrm{i} x \xi - t \lvert \xi \rvert}\, \mathrm{d}\xi = \frac{t}{\pi (t^2 + x^2)} = P_1(x, t).
\]
\end{theorem}

\begin{proof}
\uses{def:cauchy_heat_kernel, def:fourier_inv_symbol, lem:integral_ioi_cexp_neg_mul}
\leanok
Splitting $\mathbb{R} = (-\infty, 0) \cup \{x \ge 0\}$ and changing variables $\xi \mapsto -\xi$ on $(-\infty, 0)$ expresses the integral as $\int_0^\infty e^{-(t - \mathrm{i} x)\xi}\, \mathrm{d}\xi + \int_0^\infty e^{-(t + \mathrm{i} x)\xi}\, \mathrm{d}\xi$. Applying Lemma~\ref{lem:integral_ioi_cexp_neg_mul} with $\operatorname{Re}(t \mp \mathrm{i} x) = t > 0$ gives
\[
  \frac{1}{2\pi}\Bigl(\frac{1}{t - \mathrm{i} x} + \frac{1}{t + \mathrm{i} x}\Bigr) = \frac{t}{\pi (t^2 + x^2)}.
\]
\end{proof}

\begin{theorem}\label{thm:fourier_inv_gaussian}
\lean{FracLap.fourierInvFracSemigroupSymbol_two}
\uses{def:gaussian_heat_kernel, def:fourier_inv_symbol}
\leanok
For every $t > 0$ and $x \in \mathbb{R}$, the inverse Fourier transform of $\xi \mapsto e^{-t \xi^2}$ equals the Gaussian heat kernel:
\[
  \frac{1}{2\pi} \int_{\mathbb{R}} e^{\mathrm{i} x \xi - t \xi^2}\, \mathrm{d}\xi = \frac{1}{\sqrt{4 \pi t}} \exp\Bigl(-\frac{x^2}{4t}\Bigr) = P_2(x, t).
\]
\end{theorem}

\begin{proof}
\uses{def:gaussian_heat_kernel, def:fourier_inv_symbol, thm:std_gaussian_fourier}
\leanok
Because $\lvert \xi \rvert^2 = \xi^2$ and $t > 0$, Theorem~\ref{thm:std_gaussian_fourier} gives $\int_{\mathbb{R}} e^{\mathrm{i} x \xi} e^{-t \xi^2}\, \mathrm{d}\xi = (\pi / t)^{1/2} \exp(-x^2 / (4t))$. Multiplying by $(2\pi)^{-1}$ yields $\frac{1}{2\pi}\sqrt{\frac{\pi}{t}} = \frac{1}{\sqrt{4\pi t}}$, matching $P_2(x, t)$.
\end{proof}
